\section*{Capítulo \ref{ch:b1:multivar} --- Funções de Duas Variáveis}

\begin{solution}{exo:b1:multivar:1}
$\dfrac{\partial f}{\partial x} = 2xy + y\,\eu^{xy}$,
$\dfrac{\partial f}{\partial y} = x^2 + x\,\eu^{xy}$.

$\dfrac{\partial g}{\partial x} = \dfrac{2x}{x^2+y^2}$,
$\dfrac{\partial g}{\partial y} = \dfrac{2y}{x^2+y^2}$.

$\dfrac{\partial h}{\partial x} = \dfrac{-y/x^2}{1 + y^2/x^2} =
\dfrac{-y}{x^2+y^2}$,
$\dfrac{\partial h}{\partial y} = \dfrac{x}{x^2+y^2}$.
\end{solution}

\begin{solution}{exo:b1:multivar:2}
Em \omterm{prop:b1:vspaces:coordinates}{coordenadas} polares ($\rho \to 0$):

$\abs{f} = \dfrac{\rho^3\abs{\cos^2\theta\sin\theta}}{\rho^2} \leq
\rho \to 0$: \omterm{def:b1:continuity:continuous}{contínua}.

$g = \cos\theta\sin\theta$: independente de $\rho$, tomando valores diferentes ao longo de
raios diferentes (cf.\ o \cref{ex:b1:multivar:trap}): nenhum limite, não \omterm{def:b1:continuity:continuous}{contínua}.

$\abs h \leq \dfrac{\rho^3(\abs{\cos^3\theta} +
\abs{\sin^3\theta})}{\rho^2} \leq 2\rho \to 0$: \omterm{def:b1:continuity:continuous}{contínua}.
\end{solution}

\begin{solution}{exo:b1:multivar:3}
Ao longo de $y = mx$: $f(x, mx) = \dfrac{m x^3}{x^4 + m^2 x^2} =
\dfrac{mx}{x^2 + m^2} \to 0$ (para $m \neq 0$; ao longo de $y = 0$ e do
eixo $y$, $f = 0$). Logo todo limite ao longo de uma reta é $0$. Mas
na parábola $y = x^2$:
\[
f(x, x^2) = \frac{x^4}{x^4 + x^4} = \frac12 :
\]
a sequência $\bigl(\frac1n, \frac{1}{n^2}\bigr) \to (0,0)$ tem $f \to \frac12 \neq 0$. Não \omterm{def:b1:continuity:continuous}{contínua}: as retas não bastam para
testar limites de duas variáveis.
\end{solution}

\begin{solution}{exo:b1:multivar:4}
$\frac{\partial f}{\partial x} = 3x^2y^2 + y\cos(xy)$; então
\[
\frac{\partial^2 f}{\partial y\,\partial x}
= 6x^2 y + \cos(xy) - xy\sin(xy) .
\]
$\frac{\partial f}{\partial y} = 2x^3 y + x\cos(xy)$; então
\[
\frac{\partial^2 f}{\partial x\,\partial y}
= 6x^2 y + \cos(xy) - xy\sin(xy) :
\]
iguais, como Schwarz promete.
\end{solution}

\begin{solution}{exo:b1:multivar:5}
Pelo \cref{thm:b1:multivar:chain} com $(x, y) = (\cos t, \sin t)$:
\[
g'(t) = -\sin t\,\frac{\partial f}{\partial x}(\cos t, \sin t)
+ \cos t\,\frac{\partial f}{\partial y}(\cos t, \sin t)
= \bigl\langle \nabla f,\ (-\sin t, \cos t)\bigr\rangle .
\]
Num extremo de $g$, $g'(t) = 0$: $\nabla f$ é \omterm{def:b1:euclid:orthogonal}{ortogonal} ao vetor tangente $(-\sin t, \cos t)$
do círculo, logo paralelo ao vetor radial $(\cos t, \sin t)$ (no plano, o
\omterm{def:b1:euclid:orthogonal}{complemento ortogonal} de um vetor unitário é a reta que ele gera,
tomada perpendicularmente). Esta é a instância mais simples de um
multiplicador de Lagrange.
\end{solution}

\begin{solution}{exo:b1:multivar:6}
$f$: $\nabla f = (2x + y - 3,\; x + 2y) = 0$: $y = -\frac x2$ e $2x - \frac x2 = 3$: $x = 2$, $y = -1$. Segunda ordem: $r = 2$,
$s = 1$, $t = 2$: $rt - s^2 = 3 > 0$, $r > 0$: mínimo local (e de fato global ---
quadrática) em $(2, -1)$, valor $f(2,-1) = -3$.

$g$: $\nabla g = (2x,\; -2y + 4) = 0$: ponto $(0, 2)$; $r = 2$, $s = 0$, $t = -2$: $rt - s^2 = -4 < 0$: sela.

$h$: $\nabla h = \bigl(4x^3 - 4(x - y),\; 4y^3 + 4(x-y)\bigr) = 0$. Somando: $x^3 + y^3 = 0$, logo $y = -x$; substituindo: $4x^3 - 8x = 0$: $x \in \{0, \pm\sqrt2\}$.
\omterm{thm:b1:multivar:critical}{Pontos críticos}: $(0,0)$, $(\sqrt2, -\sqrt2)$, $(-\sqrt2, \sqrt2)$. Em $(\pm\sqrt2, \mp\sqrt2)$: $r = 12x^2 - 4 = 20$, $s = 4$,
$t = 20$: $rt - s^2 > 0$, $r > 0$: mínimos locais (valor $h = 4 + 4 - 16 =
-8$). Em $(0,0)$:
$r = t = -4$, $s = 4$: $rt - s^2 = 0$: mudo. Examine: $h(x, x) = 2x^4 > 0$ e $h(x, -x) = 2x^4 - 8x^2 < 0$ para $x \neq 0$
pequeno: os dois sinais em toda \omterm{def:b1:topology:open}{vizinhança} --- nenhum extremo na
origem.
\end{solution}

\begin{solution}{exo:b1:multivar:7}
$x^2 - y^2$: parciais segundas $2$ e $-2$: soma $0$. $xy$: as duas parciais
segundas puras se anulam. $\eu^x\cos y$: $\frac{\partial^2}{\partial
x^2} = \eu^x\cos y$, $\frac{\partial^2}{\partial y^2} =
-\eu^x\cos y$: soma $0$. $\ln(x^2+y^2)$: pelo
\cref{exo:b1:multivar:1},
\[
\frac{\partial^2}{\partial x^2}\ln(x^2+y^2)
= \frac{2(x^2+y^2) - 4x^2}{(x^2+y^2)^2}
= \frac{2(y^2 - x^2)}{(x^2+y^2)^2},
\]
e a versão em $y$ é a sua oposta: soma $0$.
\end{solution}

\begin{solution}{exo:b1:multivar:8}
\emph{\omterm{def:b1:linmaps:projection}{Projeção}:} o plano $P$ tem normal $n = (1, 2, -1)$ e passa por $A = (4, 0, 0)$. O
ponto mais próximo da origem é $O$ projetado: $p = O + \frac{\langle A - O, n\rangle}{\norm
n^2}\,n = \frac{4}{6}(1,2,-1) = \bigl(\frac23, \frac43,
-\frac23\bigr)$, à distância
$\frac{4}{\sqrt 6}$. (Fórmula: a distância da origem ao plano $\langle X, n \rangle = c$ é $\frac{\abs c}{\norm n}$, com
$c = 4$.)

\emph{Minimização:} $f(x,y) = x^2 + y^2 + (x + 2y - 4)^2$. Escrevendo $w = x + 2y - 4$ para o último fator:
\[
\nabla f = \bigl(2x + 2w,\; 2y + 4w\bigr) = 0
\iff x = -w \text{ e } y = -2w .
\]
Substituindo na definição de $w$: $w = -w - 4w - 4$, logo $w =
-\frac23$, o que dá $x = \frac23$,
$y = \frac43$ e $z = w =
-\frac23$. O mesmo ponto obtido pela \omterm{def:b1:linmaps:projection}{projeção}, à distância $\sqrt{\frac49 + \frac{16}{9} + \frac49} = \frac{2\sqrt 6}{3} =
\frac{4}{\sqrt6}$;
é um mínimo, pois $f \to \infty$ no infinito (uma quadrática definida positiva
mais termos \omterm{def:b1:linmaps:def}{lineares}).
\end{solution}

\begin{solution}{exo:b1:multivar:9}
Com $\rho = \sqrt{x^2+y^2}$: $\frac{\partial \rho}{\partial x} =
\frac{x}{\rho}$, logo
\[
\frac{\partial f}{\partial x} = \varphi'(\rho)\,\frac{x}{\rho},
\qquad
\frac{\partial^2 f}{\partial x^2}
= \varphi''(\rho)\,\frac{x^2}{\rho^2}
+ \varphi'(\rho)\,\Bigl(\frac{1}{\rho} -
\frac{x^2}{\rho^3}\Bigr),
\]
(regras do quociente e da cadeia). Somando a expressão simétrica em
$y$: os termos em $\varphi''$ juntam $\frac{x^2 + y^2}{\rho^2} = 1$, e os termos em $\varphi'$ juntam
$\frac{2}{\rho} - \frac{x^2 +
y^2}{\rho^3} = \frac{1}{\rho}$:
\[
\Delta f = \varphi'' + \frac{\varphi'}{\rho} .
\]
Harmônicas radiais: resolva $\varphi'' + \frac{\varphi'}{\rho} = 0$: a equação de primeira ordem $u' + \frac u\rho = 0$
para $u = \varphi'$ dá $u = \frac{c}{\rho}$ (o \cref{thm:b1:diffeq:homogeneous1}), e depois $\varphi = c\ln\rho + d$. As funções harmônicas
radiais são $c\,\ln\sqrt{x^2 + y^2} + d$ --- o potencial logarítmico do \cref{exo:b1:multivar:7} e as
constantes, nada mais.
\end{solution}

\begin{solution}{exo:b1:multivar:10}
Para $z = xy$ em $(1,1)$: parciais $y = 1$ e $x = 1$, plano tangente
$z = 1 + (x - 1) + (y - 1) = x + y - 1$. Plano tangente horizontal significa que as duas parciais se
anulam, isto é, $\nabla f = 0$: pelo \cref{ex:b1:multivar:extrema}, exatamente os \omterm{thm:b1:multivar:critical}{pontos críticos}
$(0, 0)$ e $(1, 1)$, com planos horizontais $z = 0$ e $z = -1$.
``Plano tangente horizontal'' e ``\omterm{thm:b1:multivar:critical}{ponto crítico}'' são a mesma noção,
vista no gráfico e na fórmula.
\end{solution}

\begin{solution}{exo:b1:multivar:11}
\begin{enumerate}
  \item Para $y$ fixo, a função de uma variável $x \mapsto
        f(x,y)$ tem \omterm{def:b1:derivative:def}{derivada}
        nula em $\R$, logo é constante (o \cref{thm:b1:derivative:mvt}): $f(x, y) = f(0, y)$ para todo
        $x$: $f$ depende apenas de $y$.
  \item Seja $g(u, v) = f\bigl(\frac{u+v}2, \frac{u-v}2\bigr)$. Pela regra da cadeia (o \cref{thm:b1:multivar:chain}, aplicado na
        variável $v$ com $u$ congelado):
        \[
        \frac{\partial g}{\partial v}
        = \frac12\,\frac{\partial f}{\partial x}
        - \frac12\,\frac{\partial f}{\partial y} = 0 .
        \]
        Por (1), $g$ depende apenas de $u$: $g(u, v) = \varphi(u)$ com $\varphi$ de
        classe $C^1$, e invertendo a mudança de variáveis ($u = x + y$,
        $v = x - y$),
        \[
        f(x, y) = \varphi(x + y).
        \]
        Reciprocamente, toda $f$ desse tipo satisfaz $f_x = f_y =
        \varphi'$: as
        soluções são exatamente as funções de classe $C^1$ de
        $x + y$.
\end{enumerate}
\end{solution}

\begin{solution}{exo:b1:multivar:12}
No disco \omterm{def:b1:topology:open}{aberto}, um extremo seria crítico: $\nabla f = (y,
x) = 0$ apenas em $(0,0)$,
onde $f = 0$; ali há uma sela ($f(\pm
\varepsilon, \pm\varepsilon) = \varepsilon^2 > 0 > -\varepsilon^2
= f(\pm\varepsilon, \mp\varepsilon)$), logo nenhum extremo. Pelo
teorema dos valores extremos admitido, os limites são atingidos,
necessariamente no círculo de \omterm{def:b1:topology:closure}{fronteira}. Ali, com o \cref{exo:b1:multivar:5},
\[
f(\cos t, \sin t) = \cos t\sin t = \frac{\sin 2t}{2}
\in \intcc{-\tfrac12}{\tfrac12},
\]
com máximo $\frac12$ em $t = \frac\pi4, \frac{5\pi}4$ (pontos $\pm\frac{1}{\sqrt2}(1,1)$) e mínimo $-\frac12$ em $t =
\frac{3\pi}4, \frac{7\pi}4$
(pontos $\pm\frac{1}{\sqrt2}(1,-1)$). Máximo global $\frac12$, mínimo global $-\frac12$.
\end{solution}

\begin{solution}{pb:b1:multivar:1}
\textbf{1.} Desenvolvendo o lado direito:
\[
r\Bigl(h + \frac srk\Bigr)^{\!2} + \frac{rt - s^2}{r}k^2
= rh^2 + 2shk + \frac{s^2}{r}k^2 + \frac{rt - s^2}{r}k^2
= q(h,k) .
\]
Se $rt - s^2 > 0$: os dois quadrados carregam o sinal do fator $r$, e $q(h,k) = 0$
força $k = 0$ e depois $h = 0$: sinal estrito de $r$ fora da origem.
Se $rt - s^2 < 0$: $q(1, 0) = r$, ao passo que $q\bigl(-\frac sr, 1\bigr) = \frac{rt - s^2}{r}$ tem o sinal oposto: ocorrem os
dois sinais.

\textbf{2.} Se $r = 0 \neq t$, troque os papéis de $h$ e de $k$ (o quadrado
completado em $k$), com as mesmas conclusões; note que $rt - s^2 = -s^2 < 0$ assim
que $s \neq 0$, e de fato $q =
2shk + tk^2$ assume então os dois sinais ($k$ pequeno
diante de $h$). Se $r = t = 0$: $q = 2shk$, os dois sinais se e somente se
$s \neq 0$, isto é, se e somente se $rt - s^2 = -s^2 < 0$. Resumo: $q$ assume os dois
sinais $\iff rt -
s^2 < 0$; $q$ anula-se apenas na origem (definida) $\iff rt -
s^2 > 0$ (o que
força $r \neq 0$, pois $r = 0$ dá $q(1,0) =
0$).

\textbf{3.} Com $f_x = rx + sy + \beta$ e $f_y = sx + ty +
\gamma$, o desenvolvimento direto dá
\[
f(a + h, b + k) = f(a, b) + h\,f_x(a,b) + k\,f_y(a,b)
+ \tfrac12 q(h, k),
\]
sem termos de ordem superior (a função é um \omterm{def:b1:poly:def}{polinômio} de grau $2$).
Num \omterm{thm:b1:multivar:critical}{ponto crítico} a parte \omterm{def:b1:linmaps:def}{linear} se anula: $f(X_0 + H)
- f(X_0) = \frac12 q(H)$ \emph{exatamente},
de modo que o estudo de sinal das questões 1--2 classifica: mínimo
global estrito ($rt - s^2 >
0$, $r > 0$), máximo global estrito ($rt - s^2 > 0$,
$r < 0$), e sela --- os dois sinais em toda \omterm{def:b1:topology:open}{vizinhança} --- quando
$rt - s^2 <
0$. Isto demonstra o \cref{met:b1:multivar:monge} para funções quadráticas.

\textbf{4.} A forma $q - c(h^2 + k^2)$ tem dados $r - c$, $s$, $t - c$. Com $c = \frac{rt - s^2}{r + t}$
(note que $t > 0$, pois $rt > s^2 \geq 0$ e $r > 0$):
\[
(r - c)(t - c) - s^2 = rt - s^2 - c(r + t) + c^2 = c^2 \geq 0,
\]
e $r - c \geq 0$ ($c \leq r \iff rt - s^2 \leq r^2 + rt$, verdadeiro). Se $r - c > 0$, o quadrado completado da
questão 1 mostra que $q - c(h^2 + k^2) \geq 0$; se $r - c = 0$, então $s = 0$ (de $-s^2 = -c^2 \leq \dots$, que
força a quantidade exibida a ser $\geq 0$ com primeiro fator $0$) e a
forma é $(t - c)k^2 \geq
0$. Nos dois casos, $q(h,k) \geq c(h^2 + k^2)$.

\textbf{5.} Pelas questões 3--4, $f(X) = f(X_0) + \frac12 q(X -
X_0) \geq f(X_0) + \frac c2\norm{X - X_0}^2 \to +\infty$ quando $\norm X \to \infty$: $f$ é coerciva,
e a desigualdade é estrita para $X \neq X_0$: o \omterm{thm:b1:multivar:critical}{ponto crítico} é o
minimizador global único. Para a cúbica $f = x^3 + y^3 - 3xy$ do \cref{ex:b1:multivar:extrema}, nenhuma
conclusão desse tipo vale: $f(x,0)
= x^3 \to -\infty$, e o mínimo local em $(1,1)$ não é
global --- o que falhou foi a exatidão da expansão quadrática.

\textbf{6.} Desenvolvendo a norma ao quadrado:
\[
f(u,v) = u^2\norm{C_1}^2 + 2uv\,\langle C_1, C_2\rangle +
v^2\norm{C_2}^2 - 2u\,\langle C_1, b\rangle - 2v\,\langle C_2,
b\rangle + \norm b^2 ,
\]
uma função quadrática de $(u, v)$ cuja parte quadrática é $\frac12(ru^2 + 2suv + tv^2)$, com
$r = 2\norm{C_1}^2$, $s =
2\langle C_1, C_2\rangle$, $t = 2\norm{C_2}^2$.

\textbf{7.} $rt - s^2 = 4\bigl(\norm{C_1}^2\norm{C_2}^2 -
\langle C_1, C_2\rangle^2\bigr) \geq 0$ por Cauchy--Schwarz (o \cref{thm:b1:euclid:cs}), com igualdade
exatamente quando $C_1, C_2$ são proporcionais (ou um deles é nulo), isto
é, quando o par é linearmente dependente. Logo $rt - s^2 > 0 \iff (C_1, C_2)$ \omterm{def:b1:vspaces:free}{livre}.

\textbf{8.} $\nabla f = 0$ lê-se
\[
\norm{C_1}^2 u + \langle C_1, C_2\rangle v = \langle C_1,
b\rangle,
\qquad
\langle C_1, C_2\rangle u + \norm{C_2}^2 v = \langle C_2,
b\rangle,
\]
as \omterm{pb:b1:multivar:1}{equações normais}, com a matriz de Gram à esquerda; elas dizem
$\langle uC_1 + vC_2 - b,\ C_i\rangle = 0$ para $i = 1, 2$, isto é, $b - p \perp \operatorname{Vect}(C_1, C_2)$ para $p = uC_1 +
vC_2$. Pela Parte I (questões 3
e 5), o único \omterm{thm:b1:multivar:critical}{ponto crítico} é o minimizador global único, e pelo
\cref{thm:b1:euclid:projection} a caracterização ``$p \in F$, $b - p \perp F$'' identifica $p$ como a
\omterm{thm:b1:euclid:projection}{projeção ortogonal} de $b$ sobre $F = \operatorname{Vect}(C_1, C_2)$.

\textbf{9.} $\norm{b - p}^2 = \norm b^2 - 2\langle b, p\rangle +
\norm p^2$, e $\langle p, b - p\rangle = 0$ dá $\norm p^2 =
\langle p, b\rangle$: o mínimo é igual a $\norm b^2 - \langle p,
b\rangle$.
Pitágoras: $\norm b^2 = \norm p^2 + \norm{b - p}^2$ --- o vetor de dados separa-se na sua parte
explicada $p$ e na sua parte residual $b - p$, \omterm{def:b1:euclid:orthogonal}{ortogonais} entre si.

\textbf{10.} $E(\alpha, \beta) = \norm{\alpha C_1 + \beta C_2 -
b}^2$ com os $C_1, C_2, b$ indicados; as \omterm{pb:b1:multivar:1}{equações normais} da
questão 8 são, entrada a entrada,
\[
n\alpha + \Bigl(\sum x_i\Bigr)\beta = \sum y_i,
\qquad
\Bigl(\sum x_i\Bigr)\alpha + \Bigl(\sum x_i^2\Bigr)\beta = \sum
x_i y_i .
\]

\textbf{11.} Dividindo por $n$: $\alpha + \beta\bar x = \bar y$ e $\alpha\bar x + \beta\bigl(v_x + \bar x^2\bigr) = c_{xy} +
\bar x\bar y$. Substituindo $\alpha = \bar y - \beta\bar x$ na
segunda: $\beta v_x = c_{xy}$, logo
\[
\beta = \frac{c_{xy}}{v_x}, \qquad \alpha = \bar y - \beta\bar
x ,
\]
e a primeira equação diz precisamente que $(\bar x, \bar y)$ está sobre a reta.

\textbf{12.} $v_x = \frac1n\sum(x_i - \bar x)^2 \geq 0$, nulo se e somente se todo $x_i$ é igual a
$\bar x$. E $rt - s^2 = 4\bigl(n\sum
x_i^2 - (\sum x_i)^2\bigr) = 4n^2 v_x$: a condição de liberdade da questão 7 é $v_x > 0$,
isto é, os $x_i$ não todos iguais.

\textbf{13.} Com $\alpha = \bar y - \beta\bar x$, centre os dados ($\tilde x_i = x_i - \bar x$, $\tilde y_i = y_i - \bar y$):
\[
E(\alpha, \beta) = \sum_i(\tilde y_i - \beta\tilde x_i)^2
= n\,v_y - 2\beta\,n\,c_{xy} + \beta^2 n\,v_x ,
\]
minimizado em $\beta = c_{xy}/v_x$, com valor $E_{\min} =
n\bigl(v_y - \frac{c_{xy}^2}{v_x}\bigr) = n v_y(1 - \rho^2)$. Como $E_{\min} \geq 0$ e $v_y > 0$: $\rho^2 \leq 1$,
e $\abs\rho = 1$ se e somente se $E_{\min} = 0$, isto é, se e somente se todo resíduo
se anula: os pontos estão exatamente sobre a reta.

\textbf{14.} $n = 4$: $\bar x = \frac64 = 1.5$, $\bar y = 2$, $\sum x_i^2 = 14$, logo $v_x = 3.5 - 2.25 = 1.25$, $\sum x_iy_i =
0 + 1 + 6 + 12 = 19$, logo
$c_{xy} = 4.75 - 3 = 1.75$. Logo
\[
\beta = \frac{1.75}{1.25} = 1.4,
\qquad
\alpha = 2 - 1.4\times1.5 = -0.1 :
\qquad y = 1.4\,x - 0.1 .
\]
Valores ajustados $-0.1,\ 1.3,\ 2.7,\ 4.1$; resíduos $0.1,\ -0.3,\
0.3,\ -0.1$; $E_{\min} = 0.01 + 0.09 + 0.09 + 0.01 = 0.2$ (verificação: $v_y = \frac{26}4 - 4 = 2.5$ e
$4(2.5 -
\frac{1.75^2}{1.25}) = 4(2.5 - 2.45) = 0.2$).

\textbf{15.} As duas \omterm{pb:b1:multivar:1}{equações normais} da questão 10 são exatamente
$\sum_i\varepsilon_i = 0$ e $\sum_i x_i\varepsilon_i =
0$: o vetor dos resíduos é \omterm{def:b1:euclid:orthogonal}{ortogonal} a $C_1 = (1, \dots, 1)$ e a $C_2 = (x_i)$ ---
a todo o espaço do modelo. Em particular, a reta ótima equilibra
sempre os seus erros: eles somam zero.

\textbf{16.} $\frac{\dd}{\dd\alpha}\sum(y_i - \alpha)^2 =
-2\sum(y_i - \alpha) = 0$ dá $\alpha = \bar y$ (e a \omterm{def:b1:derivative:def}{derivada} segunda $2n > 0$ faz dele o
mínimo global, sendo a função uma quadrática coerciva em uma
variável). Valor mínimo $\sum(y_i - \bar y)^2 = n v_y$: a variância é o erro quadrático
irredutível de um modelo constante.

\textbf{17.} $\frac{\dd}{\dd\beta}\sum(y_i - \beta x_i)^2 =
-2\sum x_i(y_i - \beta x_i) = 0$ dá $\beta_0 = \frac{\sum
x_iy_i}{\sum x_i^2}$. Escrevendo as duas inclinações em
quantidades centradas: $\beta_0 = \frac{c_{xy} + \bar x\bar y}{v_x + \bar
x^2}$ é igual a $\frac{c_{xy}}{v_x}$ se e somente se $c_{xy}\bar x^2 = \bar
x\bar y\,v_x$,
isto é, se e somente se $\bar x = 0$ (abscissas centradas) ou $\bar y = \beta\bar x$ --- o
que significa $\alpha = 0$: a \omterm{pb:b1:multivar:1}{reta de regressão} completa já passa pela
origem.

\textbf{18.} $g(\tau) = \norm{M - P_0}^2 - 2\tau\langle M - P_0,
w\rangle + \tau^2$ é mínima em $\tau^* = \langle M - P_0,
w\rangle$, com valor $d(M,D)^2 = \norm{M - P_0}^2 - \langle M -
P_0, w\rangle^2$. No plano,
complete $w$ numa \omterm{def:b1:vspaces:free}{base} \omterm{def:b1:euclid:orthogonal}{ortonormal} $(w, n)$ com $n = \frac{(a, b)}{\sqrt{a^2+b^2}}$ (normal unitária
de $D$): então $\norm{M - P_0}^2 = \langle M - P_0,
w\rangle^2 + \langle M - P_0, n\rangle^2$, logo $d(M, D) =
\abs{\langle M - P_0, n\rangle}$, e com $aP_{0x} + bP_{0y} =
-c$:
\[
d = \frac{\abs{a x_0 + b y_0 + c}}{\sqrt{a^2 + b^2}} .
\]

\textbf{19.} De $E_{\min}/n = v_y - \frac{c_{xy}^2}{v_x}$ e $\beta^2 v_x = \frac{c_{xy}^2}{v_x}$:
\[
v_y = \beta^2 v_x + \frac{E_{\min}}n :
\]
variância total $=$ variância ao longo da reta ajustada $+$ variância
residual. (Dividindo por $v_y$: $1 = \rho^2 + (1 - \rho^2)$, e a parcela da variância
``explicada'' pela reta é $\rho^2$.)

\textbf{20.} O espaço do modelo é $\operatorname{Vect}\bigl(C_1, C_2, C_3\bigr)$ com $C_1 = (1)$, $C_2 = (x_i)$, $C_3 = (x_i^2)$;
minimizar $\norm{aC_1 + bC_2 +
cC_3 - b}^2$ leva, exatamente como na questão 8, ao sistema de
Gram $3\times3$, cuja matriz tem entradas $\langle C_i, C_j\rangle =
\sum_k x_k^{\,i+j-2}$: a matriz de momentos
exibida. Ela é invertível se e somente se $(C_1, C_2, C_3)$ é \omterm{def:b1:vspaces:free}{livre} (o \cref{exo:b1:euclid:11});
uma relação $a + bx_k + cx_k^2 = 0$ para todo $k$ faz de cada $x_k$ uma raiz de um
mesmo \omterm{def:b1:poly:def}{polinômio} de grau $\leq 2$, o que é impossível com três valores
distintos a não ser que $a
= b = c = 0$. Compare com as matrizes de momentos
do problema de fim de semana do \cref{ch:b1:det}, em que os seus determinantes
eram quadrados de valores de Vandermonde.

\textbf{21.} Primeira: $r = 2, s = 1, t = 2$, $rt - s^2 = 3 > 0$, $r > 0$: mínimo global estrito
na origem. Segunda: $r = 2, s
= 3, t = 2$, $rt - s^2 = -5 < 0$: sela. Terceira: $r = s = t =
2$, $rt - s^2 = 0$:
degenerada --- mas $x^2 + 2xy + y^2 = (x +
y)^2 \geq 0$ anula-se em toda a reta $y = -x$: um mínimo
global (não estrito) atingido ao longo de uma reta, invisível para o
teste do determinante.

\textbf{22.} Novas somas ($n = 5$): $\bar x = 3.2$, $\bar y =
1.6$, $\sum x_iy_i = 19$, logo $c_{xy} = 3.8 - 5.12 = -1.32$,
$\sum x_i^2 = 114$, logo $v_x = 22.8 - 10.24 = 12.56$. Nova inclinação:
\[
\beta = \frac{-1.32}{12.56} \approx -0.105 :
\]
um único ponto transformou uma tendência claramente crescente ($\beta = 1.4$)
numa ligeiramente decrescente. O erro quadrático cobra um resíduo
$\varepsilon^2$, de modo que um único ponto distante --- $\abs{x_{i} - \bar x}$ grande \emph{e}
resíduo grande --- domina tanto $c_{xy}$ quanto $v_x$: os \omterm{pb:b1:multivar:1}{mínimos
quadrados} são eficientes, mas não robustos.

\textbf{23.} Ponha $W = \sum w_i$ e defina $\bar x_w =
\frac1W\sum w_ix_i$, $\bar y_w$ analogamente, $v_x^w = \frac1W\sum
w_ix_i^2 - \bar x_w^2$,
$c^w_{xy} = \frac1W\sum w_ix_iy_i - \bar
x_w\bar y_w$. A \omterm{def:b1:logic:map}{aplicação} $\langle u, v\rangle_w = \sum_i w_iu_iv_i$ é um \omterm{def:b1:euclid:def}{produto interno} em $\R^n$ ($w_i > 0$ dá
a definição positiva), logo a Parte II aplica-se palavra por palavra
e as \omterm{pb:b1:multivar:1}{equações normais} dividem-se por $W$ em
\[
\alpha + \beta\bar x_w = \bar y_w,
\qquad
\alpha\bar x_w + \beta(v^w_x + \bar x_w^2) = c^w_{xy} + \bar
x_w\bar y_w ,
\]
donde $\beta = c^w_{xy}/v^w_x$ e $\alpha = \bar y_w -
\beta\bar x_w$: as mesmas fórmulas, com toda média ponderada.

\textbf{24.} $E_{\min} = \sum\varepsilon_i^2 = 0$ força todo $\varepsilon_i = 0$: todos os pontos exatamente sobre
a reta; pela questão 13 este é o caso $\abs\rho = 1$. Teste: para $(1,1), (2,3),
(3,5)$:
$\bar x = 2$, $\bar y = 3$, $v_x = \frac{14}3 - 4 =
\frac23$, $c_{xy} = \frac{22}3 - 6 = \frac43$: $\beta = 2$, $\alpha = -1$, e de fato $y_i = 2x_i - 1$ para os
três pontos; $v_y = \frac{35}3 - 9 = \frac83$ e $\rho^2 =
\frac{(4/3)^2}{(2/3)(8/3)} = 1$.

\textbf{25.} (i) Para funções quadráticas, a expansão de ordem dois é
uma \emph{identidade}, de modo que o estudo de sinal da forma
quadrática --- pura álgebra, questões 1--2 --- classifica os \omterm{thm:b1:multivar:critical}{pontos
críticos} sem nenhum teorema de Taylor admitido. (ii) As \omterm{pb:b1:multivar:1}{equações
normais} dizem ``resíduo \omterm{def:b1:euclid:orthogonal}{ortogonal} ao espaço do modelo'', de modo que
o mínimo analítico \emph{é} a \omterm{thm:b1:euclid:projection}{projeção ortogonal} do vetor de dados: o
cálculo e a geometria euclidiana calculam o mesmo objeto. (iii) A
correlação $\rho = c_{xy}/\sqrt{v_xv_y}$ mede a parcela $\rho^2$ da variância explicada pela
reta, e Cauchy--Schwarz a limita: $\abs\rho \leq 1$, com igualdade apenas para
dados alinhados. (iv) \omterm{def:b1:vspaces:free}{Bases} e liberdade (Capítulo 18), matrizes de
Gram e de momentos e os seus determinantes (Capítulos 22--23), e a
\omterm{thm:b1:euclid:projection}{projeção ortogonal} (Capítulo 23) reapareceram todas como as peças de
trabalho de um único algoritmo. As Partes II--III desenvolvem o
\emph{método dos \omterm{pb:b1:multivar:1}{mínimos quadrados}} (Legendre, Gauss) através das
suas \emph{\omterm{pb:b1:multivar:1}{equações normais}}.
\end{solution}
